\chapter{Outlook}\label{chap:Outlook} \markright {\thechapter~Outlook}
Future work should extend the range of slab systems implemented in the configurator. In particular, the ribbed timber hollow-core system should be improved by considering different insulation products and by implementing a more refined fire verification. Post-tensioning could also be extended to ribbed concrete slabs, together with an improved two-dimensional structural handling. In addition, continuity effects for timber--concrete composite systems should be included, as proposed in the revised Eurocode~5 \cite{Eurocode5_2023}.

Further studies should investigate the sensitivity of the results to changed boundary conditions. This includes the effect of increased acoustic requirements that can be already handled, as well as higher imposed and permanent loads. Such analyses would help to identify how robust the observed system rankings are under different design assumptions.

Finally, the method should be made more accessible for practising engineers. A web-based user interface could allow users to weight the sustainability assessment criteria according to the pairwise comparison method. As a first step in this direction, a prototype website was developed that enables interactive adjustment of the criterion weights and visualises the resulting pairwise system ranking. The configurator could therefore be developed into a practical decision-making tool for early-stage design.
